\chapter{TỔNG KẾT}
\label{ch:tong_ket}

\section{Kết quả đạt được}

Đối chiếu trực tiếp với sáu mục tiêu đã nêu ở Chương~\ref{ch:gioi_thieu}:

\begin{itemize}[leftmargin=1.5cm]
\item \textbf{Mục tiêu 1 --- phân vùng class-agnostic: đã đạt.} mask mAP50 $0{,}970$ và mAP50--95
$0{,}778$ trên 312 khung hình kiểm thử chia theo cảnh.

\item \textbf{Mục tiêu 2 --- đếm và định lượng đóng góp của độ sâu: đã đạt.} Bộ lọc theo độ sâu
giảm MAE từ $0{,}695$ xuống $0{,}590$ và nâng tỉ lệ đếm đúng tuyệt đối từ $60\%$ lên $68\%$; tỉ
lệ đếm thừa giảm từ $35\%$ xuống $27\%$ trong khi tỉ lệ đếm thiếu không đổi.

\item \textbf{Mục tiêu 3 --- nhận diện từ vựng mở có từ chối: đạt một phần.} Cơ chế hoạt động
đúng thiết kế với top-1 $42{,}74\%$ và top-3 $71{,}45\%$ trên mặt nạ đã khớp, nhưng độ chính xác
đầu-cuối chỉ $13{,}17\%$ do tỉ lệ khớp mặt nạ thấp ngoài miền huấn luyện.

\item \textbf{Mục tiêu 4 --- đo kích thước và kiểm tra tính hợp lý: đã đạt.} Hệ thống xuất kích
thước theo centimet kèm kết luận hợp lý cho từng chiều, hiển thị đồng nhất ở mọi đầu ra.

\item \textbf{Mục tiêu 5 --- hiệu chỉnh tỉ lệ và kiểm chứng nghiêm ngặt: đạt về mặt kiểm chứng,
không đạt về mặt cải thiện.} Thí nghiệm kiểm soát cho kết luận \texttt{PASS}, nhưng đánh giá
để-lại-một cho thấy sai số tăng từ $10{,}143$~cm lên $12{,}598$~cm ($-24{,}2\%$). Mục tiêu kiểm
chứng đã hoàn thành trọn vẹn; điều được kiểm chứng lại là một kết quả âm tính.

\item \textbf{Mục tiêu 6 --- quan hệ không gian và đồ thị cảnh: đạt một phần.} Hệ thống xuất đầy
đủ đồ thị JSON và đạt $F_1$ vi mô $1{,}0$ trên kiểm soát tổng hợp, nhưng chưa có con số nào trên
ảnh thật vì thiếu nhãn do người gán.
\end{itemize}

\section{Hạn chế}

\begin{itemize}[leftmargin=1.5cm]
\item \textbf{Hạn chế về dữ liệu.} OCID chỉ gồm cảnh trong nhà, một loại cảm biến và tập đồ vật
hữu hạn, nên kết quả phân vùng không suy rộng ra cảnh ngoài trời hay cảm biến khác. Mẫu 200 ảnh
COCO là nhỏ và phân bố lớp rất lệch, khiến độ chính xác theo từng lớp có khoảng tin cậy rộng.
Các dải kích thước tiên nghiệm là khoảng điển hình của cả lớp, không phải nhãn thật của từng cá
thể, nên cờ báo tính hợp lý chỉ bắt được sai lệch thô.

\item \textbf{Hạn chế về thiết kế nghiên cứu.} Thí nghiệm hiệu chỉnh chạy trên 50 khung hình với
249 vật đánh giá được --- đủ để kết luận về chiều hướng nhưng chưa đủ để phân tích tách bạch
theo từng nhóm lớp vật. Nhiều khung hình chỉ có một tới ba vật đủ điều kiện bỏ phiếu, khiến cơ
chế trung vị mất phần lớn khả năng bảo vệ; đề tài chưa tách riêng ảnh hưởng của số phiếu khỏi
ảnh hưởng của sai lệch hình học tương đối. Phần quan hệ không gian chưa có nhãn thật nên chỉ có
giá trị kiểm soát phần mềm, và hai vị từ \texttt{occludes}, \texttt{partially\_occluded\_by}
hoàn toàn chưa được kiểm chứng. Cuối cùng, mọi kết luận đều là quan hệ \emph{liên hệ} trong một
lần chạy tất định, không phải kết luận nhân quả có kiểm định thống kê qua nhiều lần lặp với các
hạt giống khác nhau.

\item \textbf{Hạn chế về vật liệu và cảm biến.} Hệ thống không khôi phục được hình học bị che
khuất, và cho kết quả kém trên vật trong suốt, phản chiếu, màu tối hoặc biến dạng --- những bề
mặt mà cảm biến độ sâu trả về giá trị không hợp lệ.
\end{itemize}

\section{Hướng phát triển}

\begin{enumerate}[leftmargin=1.5cm]
\item \textbf{Nâng chất lượng nhận diện đầu-cuối --- ưu tiên cao nhất.} Với tỉ lệ khớp mặt nạ chỉ
$30{,}82\%$ ngoài miền huấn luyện, đây là nút thắt lớn nhất và cũng là nguyên nhân gốc làm nhiễu
các lá phiếu hiệu chỉnh. Hướng khả thi là huấn luyện bộ phân vùng trên tập dữ liệu đa dạng hơn,
hoặc thay bằng mô hình phân vùng tổng quát như SAM \cite{kirillov2023sam} rồi giữ nguyên tầng
nhận diện phía sau.

\item \textbf{Từ bỏ giả định một hệ số cho cả ảnh.} Kết quả âm tính chỉ ra rằng sai lệch của độ
sâu đơn mắt không thuần tuý mang tính nhân. Hai phương án thay thế: ước lượng hệ số riêng cho
từng cụm độ sâu thay vì một hệ số chung, hoặc chuyển hẳn sang dùng một vật tham chiếu đã biết
kích thước --- đường đi này đã có sẵn trong hệ thống qua tham số \texttt{-{}-ref-size-cm} và
không phụ thuộc vào chất lượng nhận diện.

\item \textbf{Gán nhãn quan hệ không gian trên ảnh thật.} Cần khoảng 100--200 cặp vật được hai
người gán độc lập, có tính hệ số đồng thuận, để thay con số kiểm soát tổng hợp hiện tại bằng một
kết quả có thể bảo vệ được. Nguyên tắc bắt buộc là gán nhãn trước rồi mới tinh chỉnh ngưỡng luật,
không bao giờ tinh chỉnh trên tập dùng để đánh giá.

\item \textbf{Mở rộng quy mô và kiểm định thống kê.} Chạy lại thí nghiệm hiệu chỉnh trên toàn bộ
OCID thay vì 50 khung hình, với nhiều hạt giống, và bổ sung khoảng tin cậy cho các chỉ số chính
để phát biểu kết luận chắc chắn hơn về mặt thống kê.
\end{enumerate}

\section{Kết luận}

Đề tài đã xây dựng và đánh giá một hệ thống hoàn chỉnh đi từ ảnh RGB-D tới một mô tả cảnh định
lượng gồm mặt nạ từng vật, số lượng, danh tính, kích thước theo centimet và quan hệ không gian.
Phân vùng class-agnostic đạt kết quả mạnh trên dữ liệu thật, và thông tin độ sâu được chứng minh
là mang giá trị bổ sung có thể đo đếm được cho bài toán đếm đối tượng --- nó loại bỏ các mặt nạ
giả mà không làm mất vật thật.

Tuy nhiên, đóng góp có giá trị nhất của đề tài lại là một kết quả âm tính. Ý tưởng dùng
tri thức tiên nghiệm về kích thước để hiệu chỉnh phép đo, khi hiện thực dưới dạng một hệ số nhân
duy nhất cho cả ảnh, đã không cải thiện được sai số trên dữ liệu giữ lại --- mặc dù bộ ước lượng
vượt qua thí nghiệm kiểm soát độ nhạy. Khoảng cách giữa hai kết luận này chính là bài học phương
pháp luận trung tâm của đề tài: một thí nghiệm kiểm soát đạt chỉ chứng minh cơ chế phản ứng đúng
với loại sai lệch mà ta cố tình tạo ra, chứ không chứng minh rằng đó là loại sai lệch chi phối
trong dữ liệu thật.

Kết quả này chỉ có thể phát hiện được nhờ thủ tục đánh giá để-lại-một đã phá vỡ tính vòng tròn
vốn có của việc dùng tri thức tiên nghiệm để vừa sửa vừa nghiệm thu. Kết quả này chỉ bộc lộ được nhờ thủ tục đánh giá loại-một-ra, vốn phá vỡ tính vòng tròn của việc dùng cùng một nguồn tri thức tiên nghiệm để vừa hiệu chỉnh vừa nghiệm thu. Nếu đánh giá theo cách thông thường, phép hiệu chỉnh sẽ luôn tỏ ra có hiệu quả và hạn chế thực sự của phương pháp đã không được phát hiện. Thứ tự công việc tiếp theo vì vậy đã rõ: cải thiện chất lượng nhận diện đầu cuối trước, sau đó mới quay lại bài toán hiệu chỉnh tỉ lệ.
